% Build: latexmk -pdf GroupXX_Poster.tex   (run from the project root; rename to GroupNN_Poster before submitting)
\documentclass[9pt]{extarticle}
\usepackage[a4paper,margin=7mm]{geometry}
\usepackage[T1]{fontenc}
\usepackage{helvet}
\renewcommand{\familydefault}{\sfdefault}
\usepackage{graphicx,booktabs,enumitem,array,xcolor}
\usepackage[most]{tcolorbox}
\graphicspath{{output/charts/}}
\setlist{nosep,leftmargin=*,itemsep=0pt}
\pagestyle{empty}\setlength{\parindent}{0pt}

\definecolor{navy}{HTML}{0B2A5B}
\newcommand{\key}[1]{\textbf{\textcolor{navy}{#1}}}
\newcommand{\block}[2]{\begin{tcolorbox}[enhanced,colback=white,colframe=navy,coltitle=white,
  fonttitle=\bfseries\small,title={#1},boxrule=0.4pt,arc=1.2mm,left=1.8mm,right=1.8mm,top=0.8mm,
  bottom=0.8mm,boxsep=0.5mm,toptitle=0.4mm,bottomtitle=0.4mm,before skip=1.5mm,after skip=0mm]#2\end{tcolorbox}}
\newcommand{\fig}[2][1]{\par\centering\includegraphics[width=#1\linewidth]{#2}\par}

\begin{document}
\begin{tcolorbox}[colback=navy,colframe=navy,coltext=white,arc=1.5mm,left=3mm,right=3mm,top=1.5mm,bottom=1.5mm]
\centering
{\bfseries\LARGE Who Gets a Recycling Point?\par}
{\normalsize Comparing collection-point types across Hong Kong's 18 districts\par}
{\footnotesize \textbf{Group XX} \,|\, DSC1001/GE1356 Introduction to Data Science \,|\, Option 2: Analyse Data Solutions\par}
{\scriptsize [Member 1, ID] · [Member 2, ID] · [Member 3, ID] · [Member 4, ID] · [Member 5, ID] · [Member 6, ID]\par}
\end{tcolorbox}

\begin{minipage}[t]{0.495\linewidth}\vspace{0pt}
\block{1. Problem and questions}{%
Hong Kong landfills \key{10,510 t} of waste a day. Households recycle only \key{22\%}
against a \key{$\approx$55\%} target, and residents say a nearby bin is their main reason
to recycle (Ma et al., 2023). \textbf{How evenly does the network reach the 18 districts,
and how do the three government solutions compare?}
\begin{enumerate}[label=RQ\arabic*.,topsep=1pt]
\item How uneven is access between districts?
\item How do solutions A, B, C compare?
\item Does access follow income, and why?
\item Does access relate to waste or recycling?
\end{enumerate}
{\scriptsize \textbf{A} public-place bins · \textbf{B} GREEN@COMMUNITY · \textbf{C} smart bins}}

\block{2. Data and method}{%
\textbf{Data:} EPD collection points (8,856 after cleaning; 5,301 public; Sep 2025); C\&SD 2025
population, income and 2021 housing mix; EPD waste statistics 2022--24; 15 reports and papers
in common units.\\
\textbf{Method:} points per 10,000 residents ($n=18$ districts); Pearson $r$ and partial
correlations; six robustness checks incl.\ leave-one-out and a permutation test. Aggregated official data only; no personal data.}

\block{3. Who looks well served depends on what you count}{%
\fig[0.78]{chart6_rank_shift.png}
Public points vary \key{3.2-fold} (Wong Tai Sin 4.8 to Wan Chai 15.2), driven by
\key{public-place bins} in parks, offices and hiking roads. Without them, Kwai Tsing,
Kwun Tong and Wong Tai Sin rise to the top.}
\block{7. Limitations}{%
\begin{itemize}
\item Counts ignore capacity and hours (a 3-hour weekly Spot = a Station)
\item One-day snapshot; $n=18$; Wan Chai is influential
\item ``Public'' is the EPD's label; no walking-distance measure
\item Associations, not causes (ecological fallacy)
\end{itemize}}
\end{minipage}\hfill
\begin{minipage}[t]{0.495\linewidth}\vspace{0pt}
\block{4. Each solution follows a different driver}{%
\fig{chart7_drivers.png}
\fig[0.72]{chart8_income_vs_driver.png}
Public-place bins thin out in dense, urban districts (A). Income patterns shrink once \key{district size} (A, B) or \key{public rental housing} (C)
is held fixed. Districts get a similar number of GREEN@COMMUNITY points, so the biggest get
the least per resident.}

\block{5. Three solutions: no single winner}{%
\fig[0.9]{chart5_solutions_compared.png}
\raggedright\footnotesize \textbf{A} widest coverage, few materials. \textbf{B} most materials and the only published cost (HK\$10,100/t), thinnest coverage. \textbf{C} reaches public estates, mostly food waste.}
\block{6. Evaluation}{%
\fig[0.68]{chart9_robustness.png}
Every result keeps its direction, but only \key{3 of 20} tests are reliable. RQ4 stays
\key{unanswered}: EPD district waste data are ``indicative only''.}
\end{minipage}

\block{8. Conclusions and impact}{%
\begin{enumerate}[topsep=1pt]
\item \textbf{Measure access by type}, not total count
\item \textbf{Scale GREEN@COMMUNITY to population}: $\approx$59 more Spots lift Yuen Long,
  Sha Tin and Tuen Mun to the median
\item \textbf{Target Stores at unmanaged older private buildings}
\item \textbf{Publish output and cost} for every solution
\end{enumerate}}


\block{References and member responsibilities}{%
\tiny \textbf{References:} EPD (2025) \emph{Waste Less collection points}, \emph{Waste statistics 2024};
C\&SD (2022, 2026) district tables; Environment Bureau (2021) \emph{Waste Blueprint 2035}; HKSAR
Government (2026) LCQ2; Chan et al.\ (2024); Ma et al.\ (2023); Guo et al.\ (2018); Zhao et al.\ (2025);
DiGiacomo et al.\ (2018); Li et al.\ (2020); Yau (2012). Full list in report.
\textbf{Responsibilities:} [Member 1: \dots] · [Member 2: \dots] · [Member 3: \dots] · [Member 4: \dots] ·
[Member 5: \dots] · [Member 6: \dots]}
\end{document}
